\section{Discussion}

\subsection{RQ A: Trustworthiness Score}
Hallucination Frequency turns out to be statistically independent from the other three metrics (Pearson r between $-0.16$ and $-0.05$ with TA, ECA, and ASC). In practice, this means a model can hallucinate fairly often while still producing accurate findings when it does flag something, and the reverse is equally true. PCA backs this up: PC1 captures 66\% of variance and loads almost entirely on TA, ECA, and ASC, while HF lands on PC2. Simply averaging all four metrics equally glosses over this two-dimensional structure.

The most defensible approach to a trustworthiness score is to treat HF on its own terms: one composite for "quality when something is found" (TA + ECA + ASC), with HF handled as a separate penalty. That said, the more pressing issue is inter-auditor agreement. ICC values stay near zero on three of the four metrics (HF $0.08$, TA $0.02$, ECA $-0.05$), and only ASC clears the significance threshold at $0.16$ (p $<0.05$). Auditors disagree substantially on individual callsites. This does not sink the metrics, but it does mean individual audit scores are noisy. Only aggregated means across multiple auditors are solid enough to cite with confidence.

\subsection{RQ B: Model Benchmarking}
Gemma3:12b comes out on top, with deepseek-r1:8b, phi4:14b, and qwen2.5-coder:7b forming a cluster just behind it that is statistically indistinguishable from one another. At the bottom, llama3.2:3b and starcoder2:7b are clearly worse than the leading group (Bonferroni-corrected p $<0.005$). The notable result here is qwen2.5-coder, the only specialist code model in the set, finishing fourth behind two general-purpose models. That outcome suggests security reasoning across a Java-JS bridge draws more on broad language understanding than on code-specific training. Parameter count also seems to matter: the 12b and 14b models dominate, while the 3b model struggles noticeably. Starcoder2's collapse across every quality dimension (TA $0.38$, ECA $0.26$, ASC $0.02$ against gemma3's $4.75$, $4.55$, $0.59$) follows from its inability to produce findings in the first place rather than from a localised reasoning weakness. It is closer to non-participation than to a measurable failure mode along any single axis.

\subsection{RQ C: Failure Taxonomy}
Context window size does not appear to be the main factor driving failures. Context length correlates with performance at $|r| < 0.08$ across all models, and if anything the relationship nudges slightly positive: richer, longer contexts weakly help rather than hurt. The same holds for bridge complexity. More exposed methods correlates with marginally better trustworthiness scores (p $\approx 0.015$), which suggests that more surface area gives models more to anchor their reasoning to.

The real failure modes are elsewhere: hallucination instability and chain reasoning collapse. Llama3.2 hallucinates in 15.5\% of audits, more than four times gemma3's $3.6\%$ rate, and this does not track with context length since hallucination spikes tend to occur in shorter contexts (mean 415 chars vs.\ 565 for normal audits). The pattern points to low-parameter models confabulating findings when the bridge interface gives them little to work with. Starcoder2 fails differently: its hallucination rate is low at the audit level, but it engages with so few callsites in the first place ($23$ findings across the entire population) that the metric reflects abstention rather than restraint. Where it does engage, its effect chain awareness sits at $0.26$ on a 10-point scale, the lowest in the lineup. It surfaces issues at the symptom level but struggles to reason through multi-step exploitation paths. These are meaningfully distinct failure modes and worth keeping separate in the taxonomy: confabulation failure (llama3.2) versus shallow reasoning failure (starcoder2).
